% !TEX root = ../../tree_notes.tex
% =====================================================================
%  Chapter 6 review -- the 15-section revision contract (chapter_end.txt).
%  Revision-only: tables, one-liners, pointers back via \thmref/\probref/
%  \secref/\coderef/\propref/\defref/\exref/\figref. No derivations.
%  Summit: LC Hard 1719 Number of Ways to Reconstruct a Tree
%  (given the ancestor-pairs, hunt a parent by degree and certify transitivity).
% =====================================================================

\chaptersection{Chapter 6 Summary}

Four sections, one inverse question: \emph{given raw data a tree left behind, rebuild the tree and
prove the rebuild is the tree}. The parent/edge view (6.1) turns pairs into adjacency and reads the
root off an in-degree histogram; the degree/code view (6.2) decides whether a degree list can be a
tree at all and counts the ones that are; the identification view (6.3) uses a union-find pre-look
to say \emph{tree, forest, or cycle} in one pass; the query view (6.4) rebuilds a hidden object one
adaptive question at a time. Every entry below is a pointer, not a re-derivation.

\begin{center}
\renewcommand{\arraystretch}{1.25}
\begin{tabular}{@{}p{0.25\linewidth} p{0.47\linewidth} p{0.20\linewidth}@{}}
\toprule
\textbf{Result} & \textbf{Statement (with conditions)} & \textbf{Where} \\
\midrule
Tree certificate & For a graph on $n$ nodes, any \textbf{two} of \{$n{-}1$ edges, connected,
  acyclic\} force the third; check two, never all three. & \thmref{c25-cert} \\
Degree sum & In every tree $\sum_v\degop(v)=2(n-1)$; the count leg of the certificate in disguise.
  & \thmref{c25-degsum} \\
Read the root & A rooted structure has exactly one node of in-degree $0$ (the root) and every other
  node of in-degree exactly $1$; two parents or two roots means \emph{not a tree}. & \propref{c25-root} \\
Degree-sequence realisability & Positive $d_1,\dots,d_n$ is tree-realisable \textbf{iff}
  $\sum d_i=2(n-1)$; the leaf-peel builds one. & \thmref{c25-degreal} \\
Pr\"ufer / count & Node $v$ appears $\degop(v)-1$ times in the length-$(n{-}2)$ code; the number of
  labelled trees with degrees $d_i$ is $(n-2)!\,/\,\prod_i(\degop(v_i)-1)!$. & \propref{c25-prufdeg} \\
Operational tree test & $n$ nodes form a tree \textbf{iff} there are $n-1$ edges \emph{and} union-find
  rejects none; $k$ rejections means $k$ surplus (cyclomatic) edges. & \corref{c25-istree} \\
Query-optimal search & Against a truthful higher/lower/equal oracle the midpoint split attains
  $\lceil\log_2(n{+}1)\rceil$ guesses, and no strategy is better in the worst case. & \thmref{c25-logn} \\
\bottomrule
\end{tabular}
\end{center}

\medskip
{\small\textbf{Algorithms from this chapter -- rebuild, then certify.}}
\begin{center}
\small
\renewcommand{\arraystretch}{1.2}
\setlength{\tabcolsep}{3pt}
\begin{tabular}{@{}p{0.19\linewidth} l l p{0.29\linewidth} p{0.155\linewidth}@{}}
\toprule
\textbf{Algorithm} & \textbf{Time} & \textbf{Memory} & \textbf{When to use (deciding phrase)} & \textbf{Code} \\
\midrule
Rebuild adj, find root, certify
  & $O(n)$ & $O(n)$
  & ``edge list / parent array given, is it a valid rooted tree?''
  & \coderef{c25-build} \\
Degree-realise by leaf-peel
  & $O(n\log n)$ & $O(n)$
  & ``build \emph{a} tree with these degrees / is this degree list a tree?''
  & \coderef{c25-degbld} \\
Union-find pre-look
  & $O((n{+}m)\,\alpha)$ & $O(n)$
  & ``is it a tree? / how many components? / delete the redundant edge''
  & \coderef{c25-uf} \\
Optimal guess (binary search)
  & $O(\log n)$ & $O(1)$
  & ``find a hidden value through a higher/lower oracle, fewest questions''
  & \coderef{c25-guess} \\
Summit: degree-greedy parent hunt
  & $O(n^3)$ & $O(n^2)$
  & ``given all ancestor pairs, how many rooted trees match?''
  & \coderef{c25-summit} \\
\bottomrule
\end{tabular}
\end{center}

\chaptersection{Master Decision Flowchart}

Post-study ladder: walk top to bottom, stop at the first branch that matches. It decides
\emph{which pile of data you were handed} and hence which rebuild-and-certify tool to run -- a
follower never trusts a count that was never a certificate.

\begin{lstlisting}[language={}, basicstyle=\ttfamily\small, numbers=none, frame=single]
  You are handed DATA about a tree. What shape is the data?
  |
  |-- an EDGE LIST or PARENT ARRAY (the pairs are explicit)
  |     +-- undirected, exactly n-1 edges -> count leg + one union-find pass   [6.3]
  |     +-- rooted / directed -> read in-degree: one 0 (root), rest exactly 1  [6.1, LC 685]
  |     +-- n edges, one on a cycle -> delete the LAST rejected edge           [6.3, LC 684]
  |     +-- raw equation / relationship pairs -> rebuild adj, DFS the ratio    [6.1, LC 399]
  |
  |-- a DEGREE SEQUENCE or CODE (only statistics survive)
  |     +-- all d_i >= 1 and sum(d) == 2(n-1)?  else NOT realisable           [6.2]
  |     +-- yes -> leaf-peel (Pruefer decode) builds one tree                  [6.2]
  |     +-- "how many trees?" -> (n-2)! / prod (d_i - 1)!  (Cayley = all d free)[6.2]
  |     +-- grid / region uniformity -> recurse into 4 quadrants (quad tree)   [6.2, LC 427]
  |
  |-- "IS IT a tree / how many components?" (cycle & connectivity ask)
  |     +-- union-find: a rejected edge = a cycle; components = final set count [6.3]
  |     +-- directed acyclicity (order exists) -> topological sort, NOT UF     [6.3, LC 210]
  |
  |-- an ORACLE (you only get to ASK)
  |     +-- higher/lower/equal -> binary search, ceil(log2(n+1)) is optimal     [6.4, LC 374]
  |     +-- a wrong guess COSTS you -> min-max k + max(left,right), interval DP  [6.4, LC 375]
  |
  +-- recursion / matrix sized by n -> a chain input is O(n) deep: explicit stack,
      or a vector<vector<bool>> sized n^2 when an adj list would do. Full DSU &
      MST -> [c17]; Pruefer/Cayley proofs -> [c32]; distance/LCA oracle -> [c24].
\end{lstlisting}

\chaptersection{Common Mistakes Checklist}

\begin{itemize}[nosep, leftmargin=1.4em]
  \item \textbf{Trusting $n-1$ edges alone.} A disconnected forest can carry $n-1$ edges. The
        certificate needs \emph{two} of the three legs -- pair the count with connectivity or with
        a cycle-free union-find pass, never the count by itself (\thmref{c25-cert}, Trap~1).
  \item \textbf{Reading in-degree on an undirected list.} Feeding both directions makes the root
        look like it has a parent and every leaf look like a non-root -- a rooted instance (LC 685)
        must scan \emph{directed} edges (\propref{c25-root}, Trap~2).
  \item \textbf{Banning the wrong edge in the two-parent repair.} When a node has two parents,
        \emph{remove the later edge first}, then run the cycle scan on the rest; returning the
        second parent without testing the cycle mis-picks the answer (\probref{c25-rc2}).
  \item \textbf{Confusing the degree \emph{sum} with a \emph{realisation}.} $\sum d=2(n-1)$ with
        all $d_i\ge1$ is necessary \emph{and} sufficient -- but a $0$ or a negative sneaking in
        breaks the peel; screen for it first (\thmref{c25-degreal}, Trap~1).
  \item \textbf{Inverting the Pr\"ufer count.} Node $v$ appears $\degop(v)-1$ times \emph{in the
        code}, and the count divides by $\prod(\degop(v)-1)!$ -- not $\degop(v)!$
        (\propref{c25-prufdeg}, Trap~2).
  \item \textbf{Running union-find on a directed acyclicity ask.} ``is a topological order
        possible?'' (LC 210) is \emph{not} the undirected cycle test; union-find ignores edge
        direction and gives the wrong verdict (\secref{c25-identify}, Trap~3).
  \item \textbf{Off-by-one and overflow in the oracle.} The interval is $[\text{lo},\text{hi}]$
        closed -- move to \code{mid-1}/\code{mid+1}; and \code{(lo+hi)/2} overflows for
        $n$ up to $2^{31}{-}1$, so use \code{lo+(hi-lo)/2} (\probref{c25-guess}, Trap~1).
\end{itemize}

\chaptersection{WA Clinic}

Verdicts are symptoms; the causes are the short list this chapter drilled. Diagnose top-down --
every leaf names the training problem.

\begin{figure}[H]
\centering
\begin{tikzpicture}[diagram,
    v/.style={boxnode, fill=primaryblue!12, draw=primaryblue, text width=22mm, align=center},
    b/.style={boxnode, fill=difforange!12, draw=difforange, text width=12mm},
    l/.style={boxnode, fill=successgreen!10, draw=successgreen, text width=58mm,
              font=\scriptsize, align=left}]
  \node[v] (verdict) {\textbf{Verdict on the judge}};
  \node[b, below left=8mm and 6mm of verdict] (wa) {\textbf{WA}};
  \node[b, below right=1mm of verdict] (tle) {\textbf{TLE}};
  \node[b, below left=8mm and 1mm of verdict, xshift=30mm] (re) {\textbf{RE}};
  \node[b, below right=8mm and 6mm of verdict] (mle) {\textbf{MLE}};
  \node[l, below right=1mm and 2mm of wa.south] (wa1)
    {count-only ``tree'' flags a disconnected forest (\corref{c25-istree}); in-degree scan on an
     undirected list invents two parents (\probref{c25-rc2}); wrong banned edge (\probref{c25-rc2});
     Pr\"ufer count divides by $\degop!$ not $(\degop-1)!$ (\propref{c25-prufdeg}); oracle interval
     off by one (\probref{c25-guess})};
  \node[l, below right=1mm and 2mm of tle.south] (tle1)
    {DFS from scratch per candidate edge instead of one union-find pass (\probref{c25-rc1});
     $O(n^3)$ transitivity scan is fine at $n\le500$ but re-DFSing the graph per pair is not
     (\probref{c25-summit})};
  \node[l, below right=1mm and 2mm of re.south] (re1)
    {recursive quad-tree / tree-from-degrees build on a degenerate $n$-deep input -- go explicit
     or raise stack (\secref{c25-degree}, Trap~3); \code{(lo+hi)/2} overflow in the oracle
     (\probref{c25-guess})};
  \node[l, below right=1mm and 2mm of mle.south] (mle1)
    {an $n\times n$ \code{bool} matrix where an adjacency list suffices -- the Summit's
     $\mathit{adj}[\cdot][\cdot]$ is $O(n^2)$ on purpose at $n\le500$, but copying it per node is
     not (\probref{c25-summit})};
  \draw[flow] (verdict) -- (wa); \draw[flow] (verdict) -- (tle);
  \draw[flow] (verdict) -- (re); \draw[flow] (verdict) -- (mle);
  \draw[flow] (wa) -- (wa1); \draw[flow] (tle) -- (tle1);
  \draw[flow] (re) -- (re1); \draw[flow] (mle) -- (mle1);
\end{tikzpicture}
\caption*{Chapter~6 verdict tree -- only the failures this chapter's problems can produce.}
\end{figure}

\medskip
\noindent Backup table (same content, terse):
\begin{center}
\small
\renewcommand{\arraystretch}{1.2}
\begin{tabular}{@{}llp{0.42\linewidth}l@{}}
\toprule
\textbf{Verdict} & \textbf{Usual cause} & \textbf{One-line fix} & \textbf{Trained in} \\
\midrule
WA  & count-only ``is a tree''   & pair $n{-}1$ edges with a UF pass        & \corref{c25-istree} \\
WA  & two-parent false alarm     & scan \emph{directed} edges only          & \probref{c25-rc2} \\
WA  & Pr\"ufer count $\degop!$    & divide by $(\degop(v)-1)!$               & \propref{c25-prufdeg} \\
WA  & oracle off-by-one          & closed $[\text{lo},\text{hi}]$, move $\pm1$ & \probref{c25-guess} \\
TLE & re-DFS per candidate edge  & one union-find sweep                      & \probref{c25-rc1} \\
RE  & dense $\mathit{adj}$ matrix   & size once at $n\le500$; explicit stack & \probref{c25-summit} \\
\bottomrule
\end{tabular}
\end{center}

\chaptersection{Quick Reference Card}

\begin{quickref}[Chapter 6 at a glance]
\textbf{One question.} The rest of the book \emph{walks} a tree someone drew for you; this chapter
answers the inverse -- \emph{what data pins a tree down, and how do I rebuild and certify it} from
that data alone? Rebuild the object, then prove it is a tree with the certificate (\thmref{c25-cert}).

\medskip
\textbf{Two invariants, always.} $m=n-1$ edges (equivalently $\sum_v\degop(v)=2(n-1)$,
\thmref{c25-degsum}) and \emph{connected} (or \emph{acyclic}) -- any two of the three certificate
legs force the third (\thmref{c25-cert}). A degree list is realisable iff every $d_i\ge1$ and the
sum is exactly $2(n-1)$ (\thmref{c25-degreal}).

\medskip
\textbf{Four shadows of the same tree $T$.} Edge list / parent array $\to$ rebuild and read the
root (\secref{c25-parent}, \exref{c25-parent-run}). Degree sequence $\to$ peel and count (\secref{c25-degree},
\exref{c25-deg-run}). Component count $\to$ union-find rejections (\secref{c25-identify},
\exref{c25-id-run}). Oracle answers $\to$ adaptive halving (\secref{c25-query}, \exref{c25-query-run}).
All four describe the \emph{identical} $T$ of \figref{c25-running}.

\medskip
\textbf{Pr\"ufer one-liner.} Node $v$ shows up $\degop(v)-1$ times in the length-$(n{-}2)$ code, so
a degree sequence is realised by $(n-2)!\,/\,\prod_v(\degop(v)-1)!$ labelled trees
(\propref{c25-prufdeg}) -- Cayley's $n^{\,n-2}$ is that sum with degrees free (\chref{c32}).
\end{quickref}

\begin{lstlisting}[caption={Chapter 6 code kit: rebuild, realise, union-find, guess}, label={lst:c25-kit}]
#include <bits/stdc++.h>
using namespace std;

// (1) Edge list / parent array -> adjacency, read the root, certify.   [6.1]
vector<vector<int>> buildFromParent(const vector<int>& par, int& root, bool& ok) {
    int n = (int)par.size() - 1; vector<vector<int>> adj(n + 1); vector<int> indeg(n + 1);
    root = 0; int zeros = 0; ok = true;
    for (int v = 1; v <= n; ++v) { int p = par[v];
        if (p == 0) { root = v; ++zeros; }                 // candidate root(s)
        else if (p >= 1 && p <= n) { adj[p].push_back(v); ++indeg[v]; }
        else ok = false; }                                 // parent out of range
    if (zeros != 1) ok = false;                            // exactly one root
    for (int v = 1; v <= n; ++v) if (v != root && indeg[v] != 1) ok = false;
    return adj;                                            // + DFS from root for connectivity
}

// (2) Degree sequence -> a realising tree by the leaf-peel (Pruefer decode). [6.2]
bool realisable(const vector<int>& deg) {                  // the WHOLE test (Thm 6.2)
    long long sum = 0;
    for (int d : deg) if (d < 1) return false; else sum += d;
    return sum == 2LL * ((int)deg.size() - 1);
}

// (3) Union-find pre-look: a rejected edge closes a cycle; count components. [6.3]
struct DSU { vector<int> p;
    DSU(int n) : p(n) { iota(p.begin(), p.end(), 0); }
    int find(int x) { return p[x] == x ? x : p[x] = find(p[x]); }
    bool uni(int a, int b) { a = find(a); b = find(b);
        if (a == b) return false; p[b] = a; return true; } };   // false => cycle

// (4) Optimal guess against a higher/lower/equal oracle -- overflow-safe.    [6.4]
int guessNumber(int n) {
    int lo = 1, hi = n;
    while (lo <= hi) {
        int mid = lo + (hi - lo) / 2;                      // NOT (lo+hi)/2 -- overflow
        int g = guess(mid);                                // -1 too high, 1 too low, 0 hit
        if (g == 0) return mid;
        if (g == -1) hi = mid - 1; else lo = mid + 1;
    }
    return -1;
}
\end{lstlisting}

\chaptersection{Write-From-Memory Drills}

Five signature skeletons, each lifted from a listing already in the chapter. Cover the key
line, reproduce it from memory, then check the answer one line below.

\begin{lstlisting}[caption={Drill 1 -- rebuild a rooted tree from its parent array (from \coderef{c25-build})}]
for (int v = 1; v <= n; ++v) { int p = par[v];
    if (p == 0) { root = v; ++zeros; }
    else if (1 <= p && p <= n) { /* FILL: ______ */; ++indeg[v]; }
    else ok = false;
}
if (/* FILL: zeros != 1 */) ok = false;                    // exactly one root
\end{lstlisting}
\noindent{\scriptsize Key: append $v$ to \code{adj[p]} for the child list; validity is
\emph{one} zero in-degree and in-degree $1$ for every other node (\propref{c25-root}).}

\begin{lstlisting}[caption={Drill 2 -- the degree-sum gate (from \coderef{c25-degbld})}]
long long sum = 0;
for (int d : deg) { if (/* FILL: d < 1 */) return {}; sum += d; }
if (/* FILL: sum != 2LL*(n-1) */) return {};               // Thm 6.2 is the WHOLE test
\end{lstlisting}
\noindent{\scriptsize Key: positivity first, then the single equation $\sum d=2(n-1)$
(\thmref{c25-degreal}); if it passes, the leaf-peel never gets stuck.}

\begin{lstlisting}[caption={Drill 3 -- the Pr\"ufer multiset (from \coderef{c25-degbld})}]
for (int i = 0; i < n; ++i)
    for (int k = 1; k < /* FILL: deg[i] */; ++k) prufer.push_back(i);   // deg(v)-1 copies
\end{lstlisting}
\noindent{\scriptsize Key: node $v$ appears $\degop(v)-1$ times, so the code has length
$\sum(\degop(v)-1)=n-2$ (\propref{c25-prufdeg}).}

\begin{lstlisting}[caption={Drill 4 -- is-it-a-tree by one union-find pass (from \coderef{c25-uf})}]
bool isTree(int n, const vector<pair<int,int>>& edges) {
    if ((int)edges.size() != /* FILL: n - 1 */) return false;         // leg (1)
    DSU d(n);
    for (auto& e : edges) if (!d.uni(e.first, e.second)) return false; // a reject = a cycle
    return true;
}
\end{lstlisting}
\noindent{\scriptsize Key: $n-1$ edges \emph{and} no rejected edge $\Rightarrow$ connected and
acyclic at once (\corref{c25-istree}).}

\begin{lstlisting}[caption={Drill 5 -- optimal play against the oracle (from \coderef{c25-guess})}]
int lo = 1, hi = n;
while (lo <= hi) {
    int mid = /* FILL: lo + (hi - lo)/2 */;
    int g = guess(mid);
    if (g == 0) return mid;
    if (g == -1) /* FILL: hi = mid - 1 */; else /* FILL: lo = mid + 1 */;
}
\end{lstlisting}
\noindent{\scriptsize Key: closed interval, so shrink to \code{mid-1}/\code{mid+1}; the midpoint
guard avoids overflow at $n=2^{31}{-}1$ (\thmref{c25-logn}, Trap~1).}

\chaptersection{Complete Problem Checklist}

Every problem this chapter offers -- 4 fully solved in-section, 4 practice, one Summit, five
mixed. Tick \emph{only} after an unaided solve; a hinted attempt counts as half. Cite by name
and number -- URLs appear once, in the solved and practice entries.

\begin{center}
\footnotesize
\renewcommand{\arraystretch}{1.25}
\begin{tabular}{@{}c p{0.17\linewidth} l p{0.20\linewidth} p{0.24\linewidth} c@{}}
\toprule
\textbf{\#} & \textbf{Problem} & \textbf{Source} & \textbf{Key technique} & \textbf{What to look for} & \textbf{Done?} \\
\midrule
\multicolumn{6}{l}{\textbf{Section 6.1 -- parent array / edge list}} \\
W6.1.1 & Redundant Connection II & LC 685 & \lv4 in-degree + union-find repair & two-parent case, ban later edge, then cycle scan & $\square$ \\
P6.1.2 & Evaluate Division & LC 399 & \lv3 rebuild graph, ratio DFS & pour equation pairs into weighted adj, walk a path & $\square$ \\
\midrule
\multicolumn{6}{l}{\textbf{Section 6.2 -- degrees, sequences, codes}} \\
W6.2.1 & Construct Quad Tree & LC 427 & \lv3 recursive uniform partition & all-equal test then split into four quadrants & $\square$ \\
P6.2.2 & Course Schedule II & LC 210 & \lv3 topological order (directed) & prerequisite pairs $\to$ an ordering, or the empty list on a cycle & $\square$ \\
\midrule
\multicolumn{6}{l}{\textbf{Section 6.3 -- is-it-a-tree / union-find}} \\
W6.3.1 & Redundant Connection & LC 684 & \lv3 union-find first reject & $n$ edges, return the last edge that closed a cycle & $\square$ \\
P6.3.2 & Number of Provinces & LC 547 & \lv2 count components & matrix $\to$ union-find, answer is the final set count & $\square$ \\
\midrule
\multicolumn{6}{l}{\textbf{Section 6.4 -- queries / interactive}} \\
W6.4.1 & Guess Number Higher or Lower & LC 374 & \lv2 binary search on oracle & overflow-safe midpoint, closed-interval shrink & $\square$ \\
P6.4.2 & Guess Number Higher or Lower II & LC 375 & \lv4 min-max interval DP & wrong guess costs $k$: $k+\max(\text{lo},\text{hi})$ & $\square$ \\
\midrule
\multicolumn{6}{l}{\textbf{Summit (chapter end)}} \\
S6 & Number of Ways to Reconstruct a Tree & LC 1719 & \lv5 degree-greedy parent hunt + transitivity & equal-degree tie $\to 2$, missing containment $\to 0$ & $\square$ \\
\bottomrule
\end{tabular}
\end{center}

\chaptersection{Recommended Study Order}

Five focused days. Day 1 fixes the certificate and the rebuild; day 2 the degree/count machinery;
day 3 the union-find identification; day 4 the query model; day 5 the Summit. Solve in this order
even where it differs from section order -- each day's tool is already in hand the day before.

\begin{center}
\small
\renewcommand{\arraystretch}{1.25}
\begin{tabular}{@{}llp{0.55\linewidth}@{}}
\toprule
\textbf{Day} & \textbf{Focus} & \textbf{Tasks} \\
\midrule
1 & Read + do 6.1 & rebuild $T$ from parent array and edge list (\exref{c25-parent-run}); LC 685; LC 399; write the root-read cold \\
2 & Read + do 6.2 & realise $T$'s degree sequence by the peel (\exref{c25-deg-run}); LC 427; LC 210; count the $180$ realisations \\
3 & Read + do 6.3 & classify $T$ and $T$+bad edge (\exref{c25-id-run}); LC 684; LC 547 \\
4 & Read + do 6.4 & halve $[1,10]$ to find $x$ (\exref{c25-query-run}); LC 374; LC 375 \\
5 & Summit + revise & LC 1719 degree-greedy hunt; walk the decision flowchart on five questions \\
\bottomrule
\end{tabular}
\end{center}

\chaptersection{Difficulty Progression}

Tiers are competency gates -- each row names the \emph{skill} the tier unlocks, not just a
rating bucket.

\begin{center}
\small
\renewcommand{\arraystretch}{1.3}
\begin{tabular}{@{}lp{0.36\linewidth}p{0.34\linewidth}@{}}
\toprule
\textbf{Tier} & \textbf{Problems} & \textbf{What passing this tier proves} \\
\midrule
$\lv2$ \ $\bigstar\bigstar$
  & LC 374, 547
  & you can drive an oracle optimally and count components with one union-find pass \\
$\lv3$ \ $\bigstar\bigstar\bigstar$
  & LC 684, 427, 399, 210
  & rebuild from raw pairs, spot the redundant edge, and partition a grid recursively \\
$\lv4$ \ $\bigstar\bigstar\bigstar\bigstar$
  & LC 685, 375
  & split on in-degree \emph{and} cycle in a directed repair; shift the split when a guess costs \\
$\lv5$ \ $\bigstar\bigstar\bigstar\bigstar\bigstar$
  & LC 1719 (Summit)
  & reconstruct from partial (ancestor) data and decide existence and multiplicity \\
\bottomrule
\end{tabular}
\end{center}

\chaptersection{Reduction Gallery}

Construction is a disguise; each line collapses a family onto one rebuild-and-certify move. No new
theory -- just where the certificate is already hiding.

\begin{itemize}[nosep, leftmargin=1.4em]
  \item ``is this a tree?'' $\longrightarrow$ $n-1$ edges plus one union-find pass with no rejection
        (\corref{c25-istree}).
  \item ``delete one edge to make it a tree'' (LC 684) $\longrightarrow$ the last edge union-find
        rejects is on the unique cycle (\probref{c25-rc1}).
  \item ``one node has two parents / a cycle too'' (LC 685) $\longrightarrow$ in-degree scan splits
        the two repair cases before the cycle test (\probref{c25-rc2}).
  \item ``build \emph{a} tree with these degrees'' (LC-ish) $\longrightarrow$ check $\sum d=2(n-1)$,
        then peel leaves off a Pr\"ufer multiset (\thmref{c25-degreal}, \coderef{c25-degbld}).
  \item ``how many trees share this degree sequence?'' $\longrightarrow$ decode count
        $(n-2)!/\prod(\degop(v)-1)!$ (\propref{c25-prufdeg}).
  \item ``is this $n\times n$ grid uniform?'' (LC 427) $\longrightarrow$ a tree built from the
        data by four-way recursion, not from edges (\probref{c25-quad}).
  \item ``can I order these prerequisites?'' (LC 210) $\longrightarrow$ a \emph{directed} acyclicity
        test (topological sort), the cousin union-find cannot answer (\secref{c25-identify}).
  \item ``group accounts sharing an email'' (LC 721) $\longrightarrow$ union-find where the keys are
        strings -- the same merge mechanic as \coderef{c25-uf}.
  \item ``are these equality/inequality constraints satisfiable?'' (LC 990) $\longrightarrow$ union
        the equals, then any inequality whose two ends are already joined is a contradiction -- a
        cycle in disguise (\corref{c25-istree}).
  \item ``rebuild an N-ary tree inside a binary one'' (LC 431) $\longrightarrow$ first-child /
        next-sibling encoding, the structural mirror of \secref{c25-parent}'s parent-array rebuild.
  \item ``how many rooted trees match these ancestor pairs?'' (LC 1719) $\longrightarrow$ degree-based
        parent hunt plus transitivity containment (\probref{c25-summit}).
\end{itemize}

\chaptersection{Mixed Problem Set}

These refuse to name their pattern; decide \emph{which rebuild-and-certify move} you need before
opening the code. A hint tells you the shape of the answer, never the tool.

\begin{mixedset}
  \practice{Course Schedule}{LeetCode 207}
    {https://leetcode.com/problems/course-schedule/}{LC Medium L3}
    {``can you finish all courses'' is ``is the prerequisite digraph acyclic'' -- but note the
    direction matters, so this is the \emph{topological} twin of \secref{c25-identify}, not the
    undirected union-find test of \corref{c25-istree}: a Kahn/DFS in-degree peel decides it.}
  \practice{Regions Cut By Slashes}{LeetCode 959}
    {https://leetcode.com/problems/regions-cut-by-slashes/}{LC Hard L4}
    {scale the grid to a cell complex and union the four neighbours of every open cell; the answer
    is the \emph{component count} of \coderef{c25-uf} after you stop merging -- the exact
    ``components = final set count'' reading of \secref{c25-identify}.}
  \practice{Accounts Merge}{LeetCode 721}
    {https://leetcode.com/problems/accounts-merge/}{LC Medium L3}
    {emails that co-occur on one account belong to one person: union-find where the universe is
    email \emph{strings}, not indices. Group by the \term{find} representative -- the same merge
    mechanic as \coderef{c25-uf}, keyed differently.}
  \practice{Satisfiability of Equality Equations}{LeetCode 990}
    {https://leetcode.com/problems/satisfiability-of-equality-equations/}{LC Medium L3}
    {union every ``$a==b$'' pair, then every ``$a!=b$'' whose ends are \emph{already} joined is a
    contradiction -- an inequality edge closing a cycle over the equalities, so it is
    \corref{c25-istree}'s rejection test in a fresh costume.}
  \practice{Encode N-ary Tree to Binary Tree}{LeetCode 431}
    {https://leetcode.com/problems/encode-n-ary-tree-to-binary-tree/}{LC Hard L4}
    {a lossless encoding \emph{is} a reconstruction: map each node to \emph{first child} and
    \emph{next sibling} pointers so the N-ary tree survives as a binary one, then rebuild it exactly
    as \secref{c25-parent} rebuilds $T$ from a parent array -- the arrow reversed.}
  \practice{Remove Max Number of Edges to Keep Graph Fully Traversable}{LeetCode 1579}
    {https://leetcode.com/problems/remove-max-number-of-edges-to-keep-graph-fully-traversable/}{LC Hard L4}
    {Type-3 edges serve both players, so claim them last: union the type-1 owners, then
    the type-2, then spend \coderef{c25-uf}'s cycle test on the shared edges -- the
    maximal-removal question is \secref{c25-identify}'s "which edges are forced" run
    backwards.}
  \practice{Lexicographically Smallest Equivalent String}{LeetCode 1061}
    {https://leetcode.com/problems/lexicographically-smallest-equivalent-string/}{LC Medium L3}
    {The equalities partition 26 letters into classes; rewrite every character as its
    class's \emph{smallest} member -- \coderef{c25-uf}'s canonical-representative
    habit, one find per letter.}
  \practice{Most Stones Removed with Same Row or Column}{LeetCode 947}
    {https://leetcode.com/problems/most-stones-removed-with-same-row-or-column/}{LC Medium L3}
    {Every removal leaves the component connected, so what survives is one stone per
    component over shared row/column coordinates: union, then count the final
    representatives -- \secref{c25-identify}'s component ledger with stones as
    vertices.}
  \practice{Possible Bipartition}{LeetCode 886}
    {https://leetcode.com/problems/possible-bipartition/}{LC Medium L3}
    {Disliked pairs must land on opposite sides: BFS two-colouring answers directly, or
    union each person with all of a rival's rivals -- a failed merge is
    \corref{c25-istree}'s rejection test wearing an odd cycle.}
  \practice{Count the Number of Complete Components}{LeetCode 2685}
    {https://leetcode.com/problems/count-the-number-of-complete-components/}{LC Medium L3}
    {Group vertices with \coderef{c25-uf}, then tally vertices and edges per component:
    complete exactly when the edges hit $v(v-1)/2$. The certify-after-building step is
    \secref{c25-degree}'s degree audit counted on edges.}
\end{mixedset}


\chaptersection{The Summit -- LC 1719 Number of Ways to Reconstruct a Tree}

\begin{notebox}[Why this is the chapter's top rope]
Sections 6.1--6.4 each rebuild or verify a tree from \emph{one clean shadow}: a full edge list, a
degree sequence, a bag of union-find answers, an oracle. The Summit hands you the faintest shadow
imaginable -- only the set of \emph{ancestor--descendant pairs}, and not even the tree -- then asks
two questions at once: does \emph{any} rooted tree realise exactly these pairs, and if so are there
\emph{several}? The whole chapter converges here: you must \emph{reconstruct} a structure from
partial data (\secref{c25-parent}), reason about \emph{degrees} to guess who can be a parent
(\secref{c25-degree}), and \emph{certify} validity -- here by the transitivity of ancestry rather
than a union-find cycle count (\secref{c25-identify}). A node's candidate parent is the
\emph{lowest-degree node still connected to it}: ancestry is transitive, so every neighbour of a
node must also be a neighbour of its parent. When two nodes can occupy the same parent slot (equal
degree) the answer is ``more than one''; when a required containment fails, ``none.''
\end{notebox}

\solvedproblem{c25-summit}{Number of Ways to Reconstruct a Tree}{LeetCode 1719}
  {https://leetcode.com/problems/number-of-ways-to-reconstruct-a-tree/}{LC Hard L5}
  {degree-sorted parent hunt + ancestor-transitivity containment test; equal-degree tie $\to$ multiple}

\begin{problemstatement}[Number of Ways to Reconstruct a Tree]
You are given \code{pairs}, a list of $\text{pairs}[i]=[x_i,y_i]$ with $x_i<y_i$ and no duplicates,
describing an unknown rooted tree on the nodes that appear. A pair $[x,y]$ is present \textbf{iff}
one of $x,y$ is an \emph{ancestor} of the other (nodes on the root-to-node path, the node itself
excluded). The tree is \emph{not} binary -- a node may have any number of children. Return
\textbf{0} if no rooted tree matches, \textbf{1} if exactly one does, and \textbf{2} if more than
one does. Two trees differ when some node has a different parent. ($1\le|\text{pairs}|\le2\cdot10^{5}$,
$1\le x_i<y_i\le500$.)
\end{problemstatement}

\begin{attemptgate}[Attempt it yourself first -- \lv{5}, give yourself 50 minutes]
A fair first attempt: notice that a node connected to \emph{every} other node must be the root (its
ancestor set is everyone), and that if you fix a root, a node's parent is \emph{some} ancestor. But
enumerating trees is exponential -- the challenge is to decide existence and multiplicity
\emph{greedily} from the degrees alone.
\peek{sort the nodes by how many pairs they appear in. Who is the most likely parent of a node $x$
-- a node of larger or smaller degree?}
\nudge{a parent must be connected to all of $x$'s neighbours too (transitivity), so look for
$x$'s parent among nodes \emph{connected to} $x$ with degree at least $\degop(x)$. The lowest such
degree is the tightest fit -- the nearest ancestor.}
\fullhint{Build the adjacency and degrees. Process nodes in \emph{increasing} degree; for node $x$
scan the higher-degree nodes for the first one $y$ that is \emph{adjacent} to $x$ -- that is
$x$'s parent candidate. If $\degop(x)=\degop(y)$ the two could swap parent slots, so mark
\emph{multiple}. Verify $y$ is adjacent to \emph{every} neighbour of $x$ (containment); if any
fails, no tree exists ($0$). If more than one node ends up with no parent, more than one root
($0$). Else return $2$ if a tie was seen, otherwise $1$. Full build and proof below.}
\end{attemptgate}

\solstep{Understanding the Problem}
The input is not edges of the tree -- it is the \emph{ancestor relation} compressed into pairs. So
$[x,y]\in\text{pairs}$ does not mean $x$ is $y$'s parent, only that one is somewhere above the
other. Two structural facts drive everything. First, ancestry is \emph{transitive}: if $p$ is
$x$'s parent, then every node related to $x$ (its ancestors, and, going down, its descendants that
reach through $x$) must also be related to $p$ -- in the pair graph, $p$'s neighbour set must
\emph{contain} $x$'s neighbour set. Second, a parent has \emph{at least} as many relations as its
child, so $\degop(p)\ge\degop(x)$. The reconstruction is therefore a greedy hunt: attach each node
to the \emph{minimum-degree} neighbour that is still at least as connected as it is.

\solstep{Key Observations}
\begin{observation}
Sorting nodes by degree means when you look at $x$, every possible parent is \emph{later} in the
list. The first later node adjacent to $x$ (smallest such degree) is the closest ancestor -- the
tightest fit, exactly as \secref{c25-degree} reads structure off degrees.
\end{observation}
\begin{observation}
If $x$ and its chosen parent $y$ have the \emph{same} degree, either could be the other's parent
without changing any pair count -- so the tree is not unique: return $2$ if a valid tree exists.
This is the multiplicity signal (\propref{c25-prufdeg} counts a kindred ``swap'' freedom).
\end{observation}
\begin{observation}
Containment is the validity test: for every neighbour $z$ of $x$, $z$ must be adjacent to $y$ (or be
$y$/$x$ itself). A failure means no rooted tree can make $y$ the parent of $x$ -- return $0$. It is
the transitivity the certificate of \thmref{c25-cert} silently assumes.
\end{observation}

\solstep{Think It Through With Me}
\begin{thinkbox}
\deadend{Treat \code{pairs} as tree edges and run a union-find.}{They are \emph{not} edges -- a pair
can join a node to a grandparent. Union-find would happily accept a triangle $\{1,2,3\}$ as a
2-edge tree and miss that two of those are the same ancestor chain. The relation is ancestry, not
adjacency.}
\deadend{Try every root and every parent assignment.}{Exponential. Degrees already order the
search: a parent has degree $\ge$ its child, so a single increasing-degree sweep fixes the candidate
for each node in one pass.}
\aha{Sort by degree; for each node find the first adjacent higher-degree node as its parent, flag an
equal-degree tie as \emph{multiple}, and verify the parent contains all the node's neighbours. Count
roots; more than one root $\to 0$, a tie $\to 2$, else $1$.}
\end{thinkbox}

\solstep{The Approach Ladder}
\begin{center}\renewcommand{\arraystretch}{1.3}
\begin{tabular}{@{}c p{0.235\linewidth} p{0.30\linewidth} l c@{}}
\toprule
\textbf{\#} & \textbf{Approach} & \textbf{Idea} & \textbf{Time} & \textbf{Verdict} \\
\midrule
1 & \lv{3} pairs-as-edges union-find & treat each pair as a tree edge        & $O(m\,\alpha)$ & \textbf{WA} \\
2 & \lv{4} root + parent enumeration & fix root, recurse on ancestor sets    & exponential    & \tle \\
3 & \lv{5} degree-sorted hunt + containment & nearest higher-degree adjacent parent & $O(n^3)$ & \ac \\
\bottomrule
\end{tabular}\end{center}
Rung 1 is the tempting trap -- the pairs are ancestry, not edges, so a cycle-free count means
nothing (Rung 1's failure is precisely why \secref{c25-identify}'s test needs the \emph{right}
edges). Rung 2 is correct but explodes. Rung 3 uses the degree order to pin one parent per node and
the containment check to certify it.

\solstep{Solution}
Build an adjacency \emph{matrix} $\mathit{adj}[\cdot][\cdot]$ and neighbour lists from \code{pairs};
self-loops $\mathit{adj}[v][v]$ are set so the containment test reads uniformly. Collect the nodes
that actually appear, sort them by increasing degree. For each node $x$ (at position $i$), scan
$j>i$ for the first node $y$ adjacent to $x$ -- that is $x$'s parent candidate. If none is found,
$x$ is a root (increment the root count). Otherwise, if $\degop(x)=\degop(y)$ flag \emph{multiple},
and verify every neighbour $z$ of $x$ is adjacent to $y$; a miss returns $0$. At the end, more than
one root returns $0$; otherwise a \emph{multiple} flag returns $2$, else $1$
(\coderef{c25-summit}).

\solstep{Proof of Correctness}
\begin{proofbox}[Nearest higher-degree parent, containment, and the tie count]
\emph{Lemma 1 (a parent has degree at least its child).} If $p$ is $x$'s parent, every node related
to $x$ is related to $p$ (transitivity of ancestry), and $p$ is additionally related to $x$, so
$\degop(p)\ge\degop(x)+1>\degop(x)$ unless $p$ and $x$ have identical relation sets, in which case
$\degop(p)=\degop(x)$. Hence scanning higher-or-equal degree positions only can find $x$'s parent.
\emph{Lemma 2 (the first adjacent one is the closest ancestor).} Among the nodes adjacent to $x$
with degree $\ge\degop(x)$, the minimum-degree one is the deepest possible ancestor -- any other
adjacent node of larger degree is an ancestor of it, not of $x$ directly. Choosing it loses no
valid tree. \emph{Lemma 3 (containment is necessary and sufficient).} $y$ can be $x$'s parent iff
$y$ is adjacent to every neighbour of $x$: \emph{necessary} by transitivity; \emph{sufficient}
because attaching $x$ under such a $y$ preserves exactly the given pairs (no pair is created -- $y$
already relates to all of $x$'s relations -- and none is lost). \emph{Multiplicity and existence.}
If some $x$ shares a degree with its parent $y$, the two are interchangeable (equal $\degop$, equal relation
sets equal), so $>1$ trees exist, flagged by \emph{multiple}; if two nodes have no parent they are
two roots, impossible in one tree ($0$); if a containment fails, no parent works ($0$). Otherwise
the greedy parents are forced, giving exactly one tree. Returning $0/2/1$ by these tests is
therefore correct.
\end{proofbox}

\solstep{Implementation}
\begin{lstlisting}[caption={Number of Ways to Reconstruct a Tree (LeetCode 1719) -- degree-sorted parent hunt + containment}, label={lst:c25-summit}]
#include <bits/stdc++.h>
using namespace std;

class Solution {
public:
    int checkWays(vector<vector<int>>& pairs) {
        int N = 501;                                       // node labels are <= 500
        vector<vector<char>> adj(N, vector<char>(N, 0));   // adjacency as a bit matrix
        vector<vector<int>> nbr(N);
        for (auto& pr : pairs) {
            int A = pr[0], B = pr[1];
            adj[A][B] = adj[B][A] = 1;
            nbr[A].push_back(B); nbr[B].push_back(A);
        }
        vector<int> active;
        for (int id = 1; id < N; ++id)
            if (!nbr[id].empty()) { active.push_back(id); adj[id][id] = 1; }  // self-loop: uniform test
        sort(active.begin(), active.end(),
             [&](int a, int b) { return nbr[a].size() < nbr[b].size(); });    // increasing degree

        bool multiple = false;
        int roots = 0;
        for (int i = 0; i < (int)active.size(); ++i) {
            int x = active[i], j = i + 1;
            while (j < (int)active.size() && !adj[x][active[j]]) ++j;         // first adjacent, >= degree
            if (j < (int)active.size()) {
                int y = active[j];
                if (nbr[x].size() == nbr[y].size()) multiple = true;          // swap -> >1 tree
                for (int z : nbr[x]) if (!adj[y][z]) return 0;                // transitivity fail -> none
            } else {
                ++roots;                                                     // x has no parent -> root
            }
        }
        if (roots > 1) return 0;                                             // two roots impossible
        return multiple ? 2 : 1;
    }
};
\end{lstlisting}

\solstep{Trace}
\begin{tracebox}[Worked example -- \code{pairs = [[1,2],[2,3],[1,3]]}]
All three nodes appear; degrees $\degop(1)=\degop(2)=\degop(3)=2$. After self-loops,
$\mathit{adj}$ is complete on $\{1,2,3\}$. Sorted by degree (tie, order $1,2,3$).
\begin{itemize}[nosep,leftmargin=1.3em]
  \item $x=1$: first adjacent later node is $2$; $\degop(1)=\degop(2)$ so \emph{multiple}$=$true;
        neighbours of $1$ are $\{2,3\}$, both adjacent to $2$ \checkmark.
  \item $x=2$: first adjacent later node is $3$; tie keeps \emph{multiple} true; neighbours
        $\{1,3\}$ both adjacent to $3$ \checkmark.
  \item $x=3$: no later node, so $3$ is the \textbf{root}; $\mathit{roots}=1$.
\end{itemize}
One root and a tie $\Rightarrow$ return $\mathbf{2}$. Indeed the three pairs are the ancestor
relation of a \emph{chain} of length $3$, and either of the two non-root orders ($1\!\to\!2\!\to\!3$
or $3\!\to\!2\!\to\!1$, plus rotations) is valid -- more than one tree. \checkmark
\end{tracebox}
\begin{tracebox}[A quick ``none'' case -- \code{pairs = [[1,2],[2,3]]}]
$\degop=[0,1,2,1]$ for nodes $1,2,3$ (node $1,3$ leaves, node $2$ centre). Sorted: $1,3,2$. $x=1$
attaches to $2$ (adjacent), no tie, neighbours of $1=\{2\}$ contained in $2$ \checkmark. $x=3$
attaches to $2$ \checkmark. $x=2$ is last $\to$ root, $\mathit{roots}=1$. No tie, no containment
break $\Rightarrow$ return $\mathbf{1}$: the unique tree is $2$ as root with children $1,3$.
\end{tracebox}

\solstep{Complexity}
\begin{complexitybox}
Time $\bigO{n^3}$ in the worst case, where $n\le500$ is the number of distinct nodes: the outer
sweep is $n$, the parent scan is $O(n)$, and the containment loop inspects up to $n$ neighbours --
about $500^3=1.25\cdot10^{8}$ elementary operations at the extreme, and $\bigO{m}$ to build the
matrix ($m\le2\cdot10^{5}$). Space $\bigO{n^2}$ for the adjacency matrix ($501^2$ bytes). The
$n\le500$ label bound is what makes a dense $\mathit{adj}[\cdot][\cdot]$ affordable here -- contrast
\secref{c25-parent}, where an edge list would blow up at $n$ large.
\end{complexitybox}

\solstep{Verification and Edge Cases}
\begin{center}\renewcommand{\arraystretch}{1.2}
\begin{tabular}{@{}p{0.30\linewidth} p{0.62\linewidth}@{}}
\toprule
\textbf{Case} & \textbf{What the code does} \\
\midrule
Star (one node in every pair) & that node has max degree, becomes the single root; leaves tie-free $\to 1$ \\
Perfect chain & every consecutive degree differs by one; no tie, one root $\to 1$ \\
Two nodes same relation set & equal-degree tie fires $\to$ \emph{multiple} $\to 2$ \\
Contradictory pairs & a containment check finds a neighbour $z$ with $\mathit{adj}[y][z]=0$ $\to 0$ \\
Two components (two roots) & \code{roots}$>1$ at the end $\to 0$ \\
\bottomrule
\end{tabular}\end{center}

\solstep{Pitfalls and Variants}
\begin{warningbox}
Three classic bugs: (i) reading \code{pairs} as tree \emph{edges} and running \corref{c25-istree}'s
$n-1$-count test -- they are \emph{ancestry} pairs, so the count is meaningless (Rung 1's WA).
(ii) Forgetting the equal-degree tie: a valid tree with a swappable parent must report $2$, not $1$.
(iii) Skipping the \code{roots}$>1$ guard -- a forest of two ``trees'' has two parentless nodes and
must be $0$. Variant -- if the problem asked to \emph{build} one tree rather than count, the same
greedy parents give it directly; and the ``distinct labels $\le500$'' bound is what lets a dense
matrix stand in for the adjacency lists of \secref{c25-parent}.
\end{warningbox}

\begin{insightbox}[Key takeaways -- three things to remember]
\begin{enumerate}[nosep, leftmargin=1.6em]
  \item \emph{Partial data still pins structure.} Ancestor pairs, though not edges, order a unique
        parent per node once you sort by degree -- the degree-reading skill of \secref{c25-degree}.
  \item \emph{Transitivity is the certificate here.} Containment ($y$ relates to all of $x$'s
        relations) plays the role the union-find rejection plays in \secref{c25-identify}.
  \item \emph{Multiplicity $\iff$ an interchangeable slot.} Equal-degree adjacent nodes can swap,
        exactly as \propref{c25-prufdeg}'s count is $>1$ whenever two degrees repeat.
\end{enumerate}
\end{insightbox}

\begin{hardeningbox}
\begin{itemize}[nosep, leftmargin=1.4em]
  \item \textbf{Larger label range}: replace the $\bigO{n^2}$ \code{char} matrix with a hash set of
        neighbours and test containment by set-subtraction -- same logic, $\bigO{m}$ memory, so the
        $n\le500$ ceiling lifts.
  \item \textbf{Return a real tree}: keep the chosen $y$ as \code{parent[x]} during the sweep; the
        greedy is already a reconstruction, closing the loop with \secref{c25-parent}.
  \item \textbf{This is counting-in-disguise}: the ``how many'' $\{0,1,2\}$ answer mirrors
        \propref{c25-prufdeg}'s exact count $(n-2)!/\prod(\degop-1)!$; the full Pr\"ufer bijection
        that turns such counts into closed forms is \chref{c32}.
\end{itemize}
\end{hardeningbox}

\chaptersection{Tool Kit Inventory}

The machine assembled so far -- cumulative through Chapter~6. Chapters~1--5's rows are carried
forward; Chapter~6's rows are appended. Pointers only, never prose.

\begin{center}
\small
\renewcommand{\arraystretch}{1.25}
\begin{tabular}{@{}p{0.24\linewidth} l c p{0.27\linewidth} p{0.16\linewidth}@{}}
\toprule
\textbf{Tool} & \textbf{Since} & \textbf{Level} & \textbf{Trigger phrase} & \textbf{Where} \\
\midrule
\multicolumn{5}{l}{\itshape Chapters 1--5 -- carried forward} \\
Adjacency list + parent-guard DFS & 1.1 & \lv1 & edge list given, not pointers & \coderef{c01-build} \\
Three-slot DFS (pre/in/post) & 3.1 & \lv1 & a node sequence in a named order & \coderef{c03-dfs} \\
Root: read parent, depth, size, height & 4.2 & \lv3 & undirected edges, need above/below & \coderef{c04-root} \\
Stamp $\tin/\tout$; $O(1)$ ancestor test & 5.1 & \lv2 & ``is $u$ above $v$?'' / DFS-order listing & \coderef{c24-stamp} \\
Flatten: subtree $=$ range $[\tin,\tin{+}\size)$ & 5.2 & \lv3 & many whole-subtree queries & \coderef{c24-flatten-build} \\
\midrule
\multicolumn{5}{l}{\itshape Chapter 6 -- new this chapter} \\
Rebuild adj, read root, certify & 6.1 & \lv3 & edge list / parent array, valid rooted tree? & \coderef{c25-build} \\
Degree-sum gate + leaf-peel realisation & 6.2 & \lv3 & ``build/count a tree from degrees'' & \coderef{c25-degbld} \\
Union-find pre-look (tree? components?) & 6.3 & \lv2 & ``is it a tree / how many provinces'' & \coderef{c25-uf} \\
Optimal guess on a 3-way oracle & 6.4 & \lv2 & ``find a hidden value, fewest questions'' & \coderef{c25-guess} \\
Degree-greedy parent hunt + containment & Summit & \lv5 & ``reconstruct from ancestor pairs'' & \coderef{c25-summit} \\
\bottomrule
\end{tabular}
\end{center}

\chaptersection{Related Topics}

\begin{itemize}[nosep, leftmargin=1.4em]
  \item \chref{c03}, \chref{c04} -- backward: rebuilding from a \emph{pair of traversals}
        (in+pre) lives in \chref{c03}; this chapter rebuilds from \emph{arbitrary} data, and the
        mechanical adjacency build / parent-guarded rooting is \secref{c04-rooting} -- we add only
        the validity and root-detection decisions.
  \item \chref{c24} -- forward/backward: 6.4 \emph{uses} the distance/ancestor oracle ($\tin/\tout$,
        $\dist$, LCA) as a black box to answer queries about a hidden tree; it never re-derives it.
  \item \chref{c28} -- forward: the full Disjoint-Set-Union structure, union by rank and path
        compression live there (its greedy use in MST is Ch.~31, \chref{c17}); \secref{c25-identify} is only the one-pass pre-look.
  \item \chref{c32} -- forward: the Pr\"ufer-code bijection and Cayley's $n^{\,n-2}$ proof; here the
        code is a \emph{teaser} that explains why a degree sequence pins a counting problem
        (\propref{c25-prufdeg}).
  \item \chref{c11}, \chref{c21} -- LCA and Mo-on-trees both \emph{assume} a rebuilt, rooted tree
        -- the object this chapter produces and certifies.
\end{itemize}

\begin{lastword}[Last Word -- Chapter 6]
\emph{``Every earlier chapter walked a tree you were \textbf{given}; this chapter asked whether a
tree can be summoned back from the \textbf{marks} it left -- and proved the summoning was honest.''}

\medskip
Chapters~3--5 taught you to traverse, store, and flatten a tree that was already drawn for you.
This chapter inverted the arrow. You were handed a pile of pairs and asked to \emph{rebuild} the
shape (\secref{c25-parent}), handed a list of degrees and asked whether \emph{any} tree wears those
numbers (\secref{c25-degree}), handed an edge stream and asked \emph{is this even a tree}
(\secref{c25-identify}), and finally handed nothing but an oracle and asked to \emph{elicit} the
answer one question at a time (\secref{c25-query}). Four shadows -- edge list, parent array, degree
sequence, distance answers -- all of the \emph{same} $T$ in \figref{c25-running}, and each enough to
rebuild it.

\medskip
Carry three reflexes forward:
\begin{enumerate}[nosep, leftmargin=1.6em]
  \item \emph{Certify with two legs, never one.} $n-1$ edges \emph{and} connected, or $n-1$ edges
        \emph{and} acyclic -- any two of the certificate force the third (\thmref{c25-cert}); the
        count alone is a trap a forest happily passes.
  \item \emph{Degrees are data.} $\sum\degop=2(n-1)$ decides realisability, $(n-2)!/\prod(\degop-1)!$
        counts the trees, and the Summit's parent hunt runs entirely on degree order
        (\thmref{c25-degreal}, \propref{c25-prufdeg}).
  \item \emph{A good question halves the world.} Against an oracle, choose the query whose answers
        leave equally small sub-cases; $\lceil\log_2\rceil$ is not a convenience, it is a
        lower bound (\thmref{c25-logn}).
\end{enumerate}

\medskip
And three habits, drilled until boring: write $T$'s parent array, degree sequence, and Pr\"ufer
code by hand and check $\sum\degop=2(n-1)$ before you trust a rebuild; keep straight the
\emph{undirected} cycle test (union-find) from the \emph{directed} one (topological sort) -- they
answer different questions and each fails the other; and never answer an oracle by a linear scan
when a midpoint halves the candidate set.

\medskip
\textbf{Next:} \chref{c05} keeps building on the rebuild instinct but installs an \emph{order} inside
the tree -- a Binary Search Tree is a tree that knows its own inorder, so search, insert, delete and
rank queries all collapse into a single root-to-leaf descent, and the degeneration of that descent is
the cliffhanger that opens the balancing chapter.

\begin{center}\textcolor{softpurple}{\rule{0.4\linewidth}{0.6pt}\quad $\blacksquare$\quad \rule{0.4\linewidth}{0.6pt}}\end{center}
\end{lastword}
